\chapter*{Lista de abreviaciones}
\addcontentsline{toc}{chapter}{Lista de abreviaciones}

% Mantener en orden alfabético. Agregar cada sigla la primera vez que se use en el texto.
\begin{description}
    \item[ASL] \emph{American Sign Language} (lengua de señas americana).
    \item[CPU] \emph{Central Processing Unit} (unidad central de procesamiento).
    \item[FedAvg] \emph{Federated Averaging} (promedio federado).
    \item[FPS] \emph{Frames per second} (cuadros por segundo).
    \item[GPU] \emph{Graphics Processing Unit} (unidad de procesamiento gráfico).
    \item[IID] Independientes e idénticamente distribuidos (referido a datos).
    \item[JSON] \emph{JavaScript Object Notation} (formato de texto para datos estructurados).
    \item[LGP] \emph{Língua Gestual Portuguesa} (lengua de señas portuguesa).
    \item[LOPO] \emph{Leave One Person Out} (validación dejando fuera a una persona).
    \item[LSA] Lengua de Señas Argentina.
    \item[LSTM] \emph{Long Short-Term Memory} (memoria de corto y largo plazo).
    \item[RGB] \emph{Red, Green, Blue} (modelo de color rojo, verde y azul de una imagen).
    \item[RNN] \emph{Recurrent Neural Network} (red neuronal recurrente).
    \item[TFLite] TensorFlow Lite.
    \item[TTS] \emph{Text-to-Speech} (síntesis de voz a partir de texto).
\end{description}
